\subsection{An exact fixed-path test for a repaired attachment}

The following statement strengthens the cycle-space criterion of
Proposition~\ref{C:attachment-cycle-test}. It concerns the repaired
path construction, with its prescribed target covered by the images of
the generating maps. Relabelling that target does not affect the
statement. In particular it applies with every admissible numerical
choice in that repaired construction.

\begin{proposition}[Proved: the kernel is generated by fixed-endpoint paths]
\label{C:fixed-path-kernel}
Let $A$ be a finite stage, or a full finite substage containing $V,W$
and every prescribed incoming generator and factorization for its
retained targets, of the repaired AKR construction. Adjoin a
single quasi-filter target $Z$ for $f,g:S\rightrightarrows T$ and write
$A'$ for the resulting category. Put $s=4k+1$, where $k\geq n$, and
write $\alpha_i:T\to Z$, $i\in\mathbb Z/s$, for its generating maps.
All old objects and arrows remain unchanged. For an object $X$ let
$G_X$ have vertices $A(V,X)$ and one directed edge $b$ from $bv_1$ to
$bv_2$ for each $b\in A(W,X)$, using $A'$ when $X=Z$. Let
$K(X)=H_1(G_X;\mathbb Z)$ and $Q=\operatorname{colim}_A K$.

Let $F\subseteq G_S$ be the equalizer subgraph: its vertices are those
$c:V\to S$ with $fc=gc$, and its edges are those $b:W\to S$ with
$fb=gb$. Then
\[
 \ker\bigl(Q\longrightarrow\operatorname{colim}_{A'}K\bigr)
 =\left\langle [(g_\#-f_\#)\sigma]\ :\
  \sigma\in C_1(G_S;\mathbb Z),\quad
  \partial\sigma=w-v,\quad v,w\in F_0\right\rangle.
 \tag{FP}\label{eq:akr-fixed-path-kernel}
\]
Only pairs $v,w$ lying in the same component of $G_S$ occur. A path may
be chosen for $\sigma$; reversing edges is represented by negative
coefficients. The displayed difference is a cycle in $G_T$, and its
class in $Q$ is independent of the chosen path. It is enough to choose
one vertex from each component of $F$ and paths joining those
representatives which lie in a common component of $G_S$.

Consequently this attachment preserves $Q$, and hence the old second
homology of the binary category of elements, if
\[
 \pi_0(F)\longrightarrow\pi_0(G_S)
 \quad\hbox{is injective}.
 \tag{FC}\label{eq:akr-fixed-components}
\]
More generally, preservation holds if and only if all classes in
\eqref{eq:akr-fixed-path-kernel} are zero in $Q$; injectivity in
\eqref{eq:akr-fixed-components} is sufficient but is not asserted to be
necessary.
\end{proposition}

\begin{proof}
We first identify the cell fibers of the attachment. Every underlying
point of every prescribed repaired target is the value of a $V$-arrow.
This holds for $V$ and $W$. It is preserved by a disjoint-sum target.
For a quasi-filter target it follows because its fixed points, private
points, and each point of every repaired path target are covered by
the generating copies: the interval formulas in
Proposition C.1 of the supplied corrected construction give a nonempty
visibility interval for every path-target point. Induction gives the
assertion. Thus the vertices of each $G_X$ may be identified with its
underlying points. No assertion about arbitrary extra, unused points
adjoined to a target is needed here.

For a point $z\in Z$, put
$I_z=\{i:z\in\alpha_i(T)\}\subseteq\mathbb Z/s$.
Property (10) makes $I_z$ geodesically convex in the odd index cycle.
Indeed, between any two indices there is a unique shortest cyclic arc
of length at most $2k\leq s-2n$. Starting at either end, (10) puts the
next index of that arc in $I_z$, and iteration puts the entire arc in
$I_z$. A nonempty proper connected subset of a cycle is an interval,
so $I_z$ is either an interval or the whole cycle.

For an edge $e:W\to Z$, its visibility set is
\[
 I_e=\{i:e=\alpha_i b\text{ for some }b:W\to T\}
     =\bigcap_{r=1}^{3}I_{ev_r}.
\]
The nontrivial inclusion follows from the established AKR routing
invariant $(c_{m+1})$: if $e(W)\subseteq\alpha_i(T)$, its three image
points, in particular at least two, lie in that image, so $e$ factors
through $\alpha_i$. Each factorization is unique because $\alpha_i$
is injective. Intersections of geodesically convex sets are
geodesically convex. Consequently each nonempty $I_e$ is also an
interval or the whole cycle.

Define the neighboring legs $\lambda_i,\mu_i:S\to T$ by
\[
 (\lambda_i,\mu_i)=
 \begin{cases}
 (g,f),&0\leq i<2k,\\
 (g,g),&i=2k,\\
 (f,g),&2k<i<4k,\\
 (f,f),&i=4k.
 \end{cases}
\]
The repaired identities say
$\alpha_i\lambda_i=\alpha_{i+1}\mu_i$.
Property (11) says that the image of this composite is exactly
$\alpha_i(T)\cap\alpha_{i+1}(T)$. Therefore the fiber graph of a
vertex $z$ has one vertex at each index of $I_z$ and one edge between
each neighboring pair of such indices. For an edge $e$, the same
assertion follows by one further application of the old routing
invariant: the unique lift $b:W\to T$ in copy $i$ has its image in
$\lambda_i(S)$ whenever $i,i+1\in I_e$, so it factors uniquely
through $\lambda_i$. Its factor $d:W\to S$ gives the neighbor lift
$\mu_i d$ in copy $i+1$. Conversely every such source edge gives this
neighbor relation. Thus, in either cell dimension, every fiber graph
is a path or a circle, including a one-vertex path.

The circular fibers are precisely the equalizing cells. Here is the
explicit verification, which also identifies their generators.
Suppose a cell lies in all copies and let $d_i$ be its source lift on
the neighboring overlap at $i$. Its lift in copy $i$ is simultaneously
$\lambda_i d_i=\mu_{i-1}d_{i-1}$, with cyclic indices. The two legs
are injective on vertices and on edges. At the upper junction these
equalities imply $d_{2k+1}=d_{2k-1}$, and successive cancellation
gives
\[
 d_{2k+j}=d_{2k-j}\qquad(1\leq j\leq2k).
\]
In particular $d_{4k}=d_0$. The lower junction then gives
$gd_0=fd_{4k}=fd_0$. Working forward through the first half, and then
through the second half, forces every $d_i=d_0$. Conversely an
equalizing cell $d$ gives a circular fiber with every $d_i=d$.

Write $C_*(G)$ for the cellular chain complex of an oriented graph,
concentrated in degrees zero and one. Define a chain map
\[
 \delta:C_*(G_S)^s\longrightarrow C_*(G_T)^s,
 \qquad
 (\delta x)_i=(\lambda_i)_\#x_i-(\mu_{i-1})_\#x_{i-1},
\]
and define
$\epsilon:C_*(G_T)^s\to C_*(G_Z)$ by
$\epsilon(y)=\sum_i(\alpha_i)_\#y_i$.
The cell-fiber graphs just computed show, integrally and degree by
degree, that
\[
 0\longrightarrow C_*(F)
 \xrightarrow{\ j\ } C_*(G_S)^s
 \xrightarrow{\ \delta\ } C_*(G_T)^s
 \xrightarrow{\ \epsilon\ } C_*(G_Z)
 \longrightarrow0
 \tag{CF}\label{eq:akr-cell-fibers}
\]
is exact, where $j$ is diagonal inclusion. For a path fiber the
incidence map has zero kernel and cokernel $\mathbb Z$; for a circle
fiber its kernel is the diagonal $\mathbb Z$ and its cokernel is
$\mathbb Z$. These descriptions use bases and hence introduce no
torsion. They commute with the boundary because all displayed maps
are graph maps.

For completeness, $F$ is the subgraph induced on its fixed vertices.
If the endpoints of $b$ are fixed, then $fb$ and $gb$ share their
first two image points. Invariant $(c_m)$ factors one through the
other. The only endomorphism of the initial object $W$ is its
identity, since the construction adds only new maximal targets and
leaves old hom-sets unchanged. Hence $fb=gb$. This induced-subgraph
description is useful for testing \eqref{eq:akr-fixed-components}, but
the chain argument only uses the equalizer definition.

Let $D_* =\operatorname{im}\delta$ be the image chain complex in
$C_*(G_T)^s$. Splitting \eqref{eq:akr-cell-fibers} into two short exact
sequences of complexes gives
\[
 \ker\bigl(K(T)^s\longrightarrow K(Z)\bigr)
     =\operatorname{im}\bigl(H_1(D_*)\longrightarrow K(T)^s\bigr),
\]
and
\[
 K(S)^s\longrightarrow H_1(D_*)\longrightarrow H_0(F)
       \xrightarrow{\ H_0(j)\ }H_0(G_S)^s.
 \tag{HC}\label{eq:akr-fiber-homology}
\]
The last map is diagonal on the component-inclusion homomorphism.
Its kernel is therefore exactly
$\ker(H_0(F)\to H_0(G_S))$. That kernel is generated by differences
$[w]-[v]$ of components of $F$ lying in one component of $G_S$.

Choose a graph path $\sigma$ from $v$ to $w$. Its diagonal copy in
$C_1(G_S)^s$ has boundary $j(w-v)$, so
$\delta(\sigma,\ldots,\sigma)$ is a cycle in $D_*$ lifting that
component difference, up to the immaterial common boundary sign.
Its $i$th component
\[
 \kappa_i=(\lambda_i)_\#\sigma-(\mu_{i-1})_\#\sigma
\]
is separately a cycle in $G_T$, because $f$ and $g$ agree at $v,w$.
The sum of these components is
\[
 \sum_i\kappa_i
 =\sum_i\bigl((\lambda_i)_\#-(\mu_i)_\#\bigr)\sigma
 =(g_\#-f_\#)\sigma.
\]
Here there are $2k$ occurrences of $(g,f)$, $2k-1$ of $(f,g)$,
and two zero differences. On the other hand, the image in $Q$ of
every element coming from $K(S)^s$ in
\eqref{eq:akr-fiber-homology} is zero: each of its terms is the
difference of the images of a source cycle along two old arrows.
Proposition~\ref{C:attachment-cycle-test} now gives exactly
\eqref{eq:akr-fixed-path-kernel}, in both directions.

If $\sigma'$ has the same boundary as $\sigma$, their difference is
a cycle $\tau\in K(S)$, and
$[(g_\#-f_\#)\tau]=[\tau]-[\tau]=0$ in $Q$. Moving either endpoint
within its component of $F$ also leaves the class unchanged, since
$f_\#=g_\#$ on chains of $F$. These observations establish all path
and representative assertions. Finally, the dimension-shifting
formula identifies the old second homology with the kernel of the
augmentation $Q\to\mathbb Z$; every generator in
\eqref{eq:akr-fixed-path-kernel} has augmentation zero. Thus the
criterion for preserving $Q$ is also the criterion for preserving
that old second homology.
\end{proof}

\paragraph{Exact remaining issue.}
Proposition~\ref{C:fixed-path-kernel} does not establish
\eqref{eq:akr-fixed-components} for arbitrary stages. The finite tests
of that condition therefore address the exact possible attachment
kernel, but passing those tests does not prove its eventual
vanishing. Conversely a disconnected fixed subgraph alone would not
prove a nonzero kernel: one must compute the associated class
$[(g_\#-f_\#)\sigma]$ in the old colimit. Even a nonzero kernel would
only show that some old second-homology class dies; it would not by
itself show that the full infinite category is total or that its
binary second homology vanishes.
